% GENERATED FILE. Edit canonical lessons, then run npm run generate:books.
% canonical-source-hash: fnv1a64:7c365c8fa228cef1
% canonical-lessons: JA-C451-shukkin-suru, JA-C451-tsuukin-suru, JA-C451-chikoku-suru, JA-C451-kyuukei-suru, JA-C451-zangyou-suru

\chapter{Off to Work}
\label{ch:ja-off-to-work}
\hlchaptermodality{451}

\begin{chapteropening}
\textbf{By the end of this chapter:} \emph{I can talk about going to work and working late.}

Complete the last lesson of chapter 451: I can talk about going to work and working late.
\end{chapteropening}

\section[\texorpdfstring{\ja{しゅっきんする} (shukkin suru)}{しゅっきんする (shukkin suru)}]{\ja{しゅっきんする} (shukkin suru) — to go in to work}

\label{lesson:JA-C451-shukkin-suru}



\begin{quote}
Before the new one: say the Japanese for eye drops, then the Japanese for to nurse the sick.
\end{quote}

\subsection*{You'll want to know: \ja{しゅっきんする}}
\textbf{\ja{しゅっきんする}} — \emph{shukkin suru} — \textquotedblleft{}to go in to work\textquotedblright{}.

\begin{quote}
\textbf{\ja{ちちは} \ja{どようびも} \ja{しゅっきんしました。}} — \emph{chichi wa doyōbi mo shukkin shimashita} — \textquotedblleft{}My father went in to work on Saturday too.\textquotedblright{}
\end{quote}

Polite \textbf{\ja{しゅっきんします}}, \emph{shukkinshimasu}; a \ja{する} verb. Past (polite): \textbf{\ja{ちちが} \ja{しゅっきんしました。}}, \emph{chichi ga shukkin shimashita}.

\begin{cousinweb}
\textbf{\ja{しゅっきんする}} (\emph{shukkin suru}) is Sino-Japanese (kango) \emph{shukkin} with the native verb \emph{suru}, \textquotedblleft{}to do\textquotedblright{}: two Chinese characters read with their on'yomi, readings borrowed from Chinese. They are \emph{shu}, \textquotedblleft{}go out\textquotedblright{} (Middle Chinese \emph{tsyhwit}) and \emph{kin}, \textquotedblleft{}serve\textquotedblright{} (\emph{gj+n}) (KANJIDIC2; Middle Chinese in Baxter's transcription).
\end{cousinweb}

\begin{grammarlens}[title={What you've built}]
One more word for describing what you lived and felt.
\end{grammarlens}

\subsection*{Your turn}
\begin{itemize}
\raggedright
  \item \emph{Say it:} \emph{shukkin suru}
  \item \emph{Say it:} \emph{shukkin suru}, once more
  \item \emph{Say it:} \emph{chichi wa doyōbi mo shukkin shimashita} — My father went in to work on Saturday too.
  \item \emph{Recall:} say the Japanese for eye drops, then the Japanese for to nurse the sick, then say \emph{shukkin suru} again
\end{itemize}

\begin{tcolorbox}[breakable,colback=teal!4,colframe=teal!35!black,title={Before you move on}]
What does \ja{しゅっきんする} mean? (\textquotedblleft{}To go in to work\textquotedblright{}.) Say it once more.
\end{tcolorbox}
\section[\texorpdfstring{\ja{つうきんする} (tsūkin suru)}{つうきんする (tsūkin suru)}]{\ja{つうきんする} (tsūkin suru) — to commute}

\label{lesson:JA-C451-tsuukin-suru}



\begin{quote}
Before the new one: say the Japanese for to nurse the sick, then the Japanese for to go in to work.
\end{quote}

\subsection*{You'll want to know: \ja{つうきんする}}
\textbf{\ja{つうきんする}} — \emph{tsūkin suru} — \textquotedblleft{}to commute\textquotedblright{}.

\begin{quote}
\textbf{\ja{あには} \ja{でんしゃで} \ja{つうきんしています。}} — \emph{ani wa densha de tsūkin shite imasu} — \textquotedblleft{}My elder brother commutes by train.\textquotedblright{}
\end{quote}

Polite \textbf{\ja{つうきんします}}, \emph{tsūkinshimasu}; a \ja{する} verb. Past (polite): \textbf{\ja{ははが} \ja{じてんしゃで} \ja{つうきんしました。}}, \emph{haha ga jitensha de tsūkin shimashita}.

\begin{cousinweb}
\textbf{\ja{つうきんする}} (\emph{tsūkin suru}) is Sino-Japanese (kango) \emph{tsūkin} with the native verb \emph{suru}, \textquotedblleft{}to do\textquotedblright{}: two Chinese characters read with their on'yomi, readings borrowed from Chinese. They are \emph{tsū}, \textquotedblleft{}commute\textquotedblright{} (Middle Chinese \emph{thuwng}) and \emph{kin}, \textquotedblleft{}serve\textquotedblright{} (\emph{gj+n}) (KANJIDIC2; Middle Chinese in Baxter's transcription).
\end{cousinweb}

\begin{grammarlens}[title={What you've built}]
One more word for describing what you lived and felt.
\end{grammarlens}

\subsection*{Your turn}
\begin{itemize}
\raggedright
  \item \emph{Say it:} \emph{tsūkin suru}
  \item \emph{Say it:} \emph{tsūkin suru}, once more
  \item \emph{Say it:} \emph{ani wa densha de tsūkin shite imasu} — My elder brother commutes by train.
  \item \emph{Recall:} say the Japanese for to nurse the sick, then the Japanese for to go in to work, then say \emph{tsūkin suru} again
\end{itemize}

\begin{tcolorbox}[breakable,colback=teal!4,colframe=teal!35!black,title={Before you move on}]
What does \ja{つうきんする} mean? (\textquotedblleft{}To commute\textquotedblright{}.) Say it once more.
\end{tcolorbox}
\section[\texorpdfstring{\ja{ちこくする} (chikoku suru)}{ちこくする (chikoku suru)}]{\ja{ちこくする} (chikoku suru) — to arrive late}

\label{lesson:JA-C451-chikoku-suru}



\begin{quote}
Before the new one: say the Japanese for to go in to work, then the Japanese for to commute.
\end{quote}

\subsection*{You'll want to know: \ja{ちこくする}}
\textbf{\ja{ちこくする}} — \emph{chikoku suru} — \textquotedblleft{}to arrive late\textquotedblright{}.

\begin{quote}
\textbf{\ja{あさねぼうして、ちこくしました。}} — \emph{asanebō shite, chikoku shimashita} — \textquotedblleft{}I overslept and arrived late.\textquotedblright{}
\end{quote}

Polite \textbf{\ja{ちこくします}}, \emph{chikokushimasu}; a \ja{する} verb. Past (polite): \textbf{\ja{がくせいが} \ja{ちこくしました。}}, \emph{gakusei ga chikoku shimashita}.

\begin{cousinweb}
\textbf{\ja{ちこくする}} (\emph{chikoku suru}) is Sino-Japanese (kango) \emph{chikoku} with the native verb \emph{suru}, \textquotedblleft{}to do\textquotedblright{}: two Chinese characters read with their on'yomi, readings borrowed from Chinese. They are \emph{chi}, \textquotedblleft{}late\textquotedblright{} (Middle Chinese \emph{drij}) and \emph{koku}, \textquotedblleft{}time\textquotedblright{} (\emph{khok}) (KANJIDIC2; Middle Chinese in Baxter's transcription).
\end{cousinweb}

\begin{grammarlens}[title={What you've built}]
One more word for describing what you lived and felt.
\end{grammarlens}

\subsection*{Your turn}
\begin{itemize}
\raggedright
  \item \emph{Say it:} \emph{chikoku suru}
  \item \emph{Say it:} \emph{chikoku suru}, once more
  \item \emph{Say it:} \emph{asanebō shite, chikoku shimashita} — I overslept and arrived late.
  \item \emph{Recall:} say the Japanese for to go in to work, then the Japanese for to commute, then say \emph{chikoku suru} again
\end{itemize}

\begin{tcolorbox}[breakable,colback=teal!4,colframe=teal!35!black,title={Before you move on}]
What does \ja{ちこくする} mean? (\textquotedblleft{}To arrive late\textquotedblright{}.) Say it once more.
\end{tcolorbox}
\section[\texorpdfstring{\ja{きゅうけいする} (kyūkei suru)}{きゅうけいする (kyūkei suru)}]{\ja{きゅうけいする} (kyūkei suru) — to take a break}

\label{lesson:JA-C451-kyuukei-suru}



\begin{quote}
Before the new one: say the Japanese for to commute, then the Japanese for to arrive late.
\end{quote}

\subsection*{You'll want to know: \ja{きゅうけいする}}
\textbf{\ja{きゅうけいする}} — \emph{kyūkei suru} — \textquotedblleft{}to take a break\textquotedblright{}.

\begin{quote}
\textbf{\ja{ひるに} \ja{すこし} \ja{きゅうけいします。}} — \emph{hiru ni sukoshi kyūkei shimasu} — \textquotedblleft{}I take a short break at noon.\textquotedblright{}
\end{quote}

Polite \textbf{\ja{きゅうけいします}}, \emph{kyūkeishimasu}; a \ja{する} verb. Past (polite): \textbf{\ja{どうりょうが} \ja{きゅうけいしました。}}, \emph{dōryō ga kyūkei shimashita}.

\begin{cousinweb}
\textbf{\ja{きゅうけいする}} (\emph{kyūkei suru}) is Sino-Japanese (kango) \emph{kyūkei} with the native verb \emph{suru}, \textquotedblleft{}to do\textquotedblright{}: two Chinese characters read with their on'yomi, readings borrowed from Chinese. They are \emph{kyū}, \textquotedblleft{}rest\textquotedblright{} (Middle Chinese \emph{xjuw}) and \emph{kei}, \textquotedblleft{}recess\textquotedblright{} (\emph{khjejH}) (KANJIDIC2; Middle Chinese in Baxter's transcription).
\end{cousinweb}

\begin{grammarlens}[title={What you've built}]
One more word for describing what you lived and felt.
\end{grammarlens}

\subsection*{Your turn}
\begin{itemize}
\raggedright
  \item \emph{Say it:} \emph{kyūkei suru}
  \item \emph{Say it:} \emph{kyūkei suru}, once more
  \item \emph{Say it:} \emph{hiru ni sukoshi kyūkei shimasu} — I take a short break at noon.
  \item \emph{Recall:} say the Japanese for to commute, then the Japanese for to arrive late, then say \emph{kyūkei suru} again
\end{itemize}

\begin{tcolorbox}[breakable,colback=teal!4,colframe=teal!35!black,title={Before you move on}]
What does \ja{きゅうけいする} mean? (\textquotedblleft{}To take a break\textquotedblright{}.) Say it once more.
\end{tcolorbox}
\section[\texorpdfstring{\ja{ざんぎょうする} (zangyō suru)}{ざんぎょうする (zangyō suru)}]{\ja{ざんぎょうする} (zangyō suru) — to work overtime}

\label{lesson:JA-C451-zangyou-suru}



\begin{quote}
Before the new one: say the Japanese for to arrive late, then the Japanese for to take a break.
\end{quote}

\subsection*{You'll want to know: \ja{ざんぎょうする}}
\textbf{\ja{ざんぎょうする}} — \emph{zangyō suru} — \textquotedblleft{}to work overtime\textquotedblright{}.

\begin{quote}
\textbf{\ja{こんばんも} \ja{ざんぎょうします。}} — \emph{konban mo zangyō shimasu} — \textquotedblleft{}I am working overtime tonight too.\textquotedblright{}
\end{quote}

Polite \textbf{\ja{ざんぎょうします}}, \emph{zangyōshimasu}; a \ja{する} verb. Past (polite): \textbf{\ja{ちちが} \ja{ざんぎょうしました。}}, \emph{chichi ga zangyō shimashita}.

\begin{cousinweb}
\textbf{\ja{ざんぎょうする}} (\emph{zangyō suru}) is Sino-Japanese (kango) \emph{zangyō} with the native verb \emph{suru}, \textquotedblleft{}to do\textquotedblright{}: two Chinese characters read with their on'yomi, readings borrowed from Chinese. They are \emph{zan}, \textquotedblleft{}leftover\textquotedblright{} (Middle Chinese \emph{dzan}) and \emph{gyō}, \textquotedblleft{}business\textquotedblright{} (\emph{ngjaep}) (KANJIDIC2; Middle Chinese in Baxter's transcription).
\end{cousinweb}

\begin{grammarlens}[title={What you've built}]
One more word for describing what you lived and felt.
\end{grammarlens}

\subsection*{Your turn}
\begin{itemize}
\raggedright
  \item \emph{Say it:} \emph{zangyō suru}
  \item \emph{Say it:} \emph{zangyō suru}, once more
  \item \emph{Say it:} \emph{konban mo zangyō shimasu} — I am working overtime tonight too.
  \item \emph{Recall:} say the Japanese for to arrive late, then the Japanese for to take a break, then say \emph{zangyō suru} again
\end{itemize}

\begin{tcolorbox}[breakable,colback=teal!4,colframe=teal!35!black,title={Before you move on}]
What does \ja{ざんぎょうする} mean? (\textquotedblleft{}To work overtime\textquotedblright{}.) Say it once more.
\end{tcolorbox}
